\begin{table}[htbp]
    \centering
    \caption{Timeline of Generic Prescribing Rate by Authority (Overall Providers)}
    \label{tab:partd_generic_rate_timeline}
    \vspace{0.2cm}
    \resizebox{\textwidth}{!}{
    \begin{tabular}{lccccccccccc}
        \toprule
        \textbf{State NP Authority} & \textbf{2013} & \textbf{2014} & \textbf{2015} & \textbf{2016} & \textbf{2017} & \textbf{2018} & \textbf{2019} & \textbf{2020} & \textbf{2021} & \textbf{2022} & \textbf{2023} \\
        \midrule
        Restricted & 0.627 & 0.646 & 0.673 & 0.690 & 0.703 & 0.719 & 0.732 & 0.737 & 0.737 & 0.736 & 0.736 \\
        Reduced & 0.631 & 0.651 & 0.673 & 0.691 & 0.709 & 0.724 & 0.734 & 0.738 & 0.739 & 0.739 & 0.739 \\
        Full Practice & 0.658 & 0.674 & 0.694 & 0.712 & 0.728 & 0.742 & 0.752 & 0.756 & 0.757 & 0.754 & 0.754 \\
        \midrule
        \textbf{National Average} & 0.638 & 0.657 & 0.680 & 0.697 & 0.713 & 0.727 & 0.739 & 0.743 & 0.744 & 0.742 & 0.742 \\
        \bottomrule
    \end{tabular}
    }
    \vspace{0.1cm}
    \begin{minipage}{\textwidth}
        \footnotesize \textit{Notes:} Rate: Proportion of total Part D prescriptions filled with generic drugs. Higher implies cost efficiency. Categorization reflects static 2023 status. Aggregated across all provider types.
    \end{minipage}
\end{table}